\section{Introduction}
\label{sec:introduction}

Large simulation, optimization, and control pipelines rarely solve one equation
instance in isolation. They repeatedly solve related sparse linear systems,
nonlinear residual equations, ordinary differential equations (ODEs),
differential--algebraic equations (DAEs), and event-driven subsystems as parameters,
states, or boundary conditions change. Mature numerical libraries provide robust
general-purpose algorithms for these families~\cite{saad1986gmres,li2005superlu,
hindmarsh2005sundials,amestoy2001mumps}. However, the best algorithm and hardware
path depend on structure, scale, conditioning, reuse, and requested accuracy. A
single globally fixed solver consequently leaves substantial performance on the
table.

Learned numerical components create a second opportunity and a second failure
mode. Neural operators and learned preconditioners can amortize computation over a
parameterized family~\cite{li2021fno,lu2021deeponet,li2023preconditioner}, but raw
prediction time is not complete solve time. Candidate construction, data movement,
local correction, original-equation verification, and the next numerical attempt after
rejecting an unsuitable candidate can dominate short solves. Moreover, a low average test error does not establish
that an individual candidate satisfies the original equation. Directly replacing
a classical solver therefore couples performance uncertainty to numerical risk.

This paper studies the following scientific question:

\begin{quote}
\emph{Can heterogeneous classical and learned solver experts be fused so that
selection minimizes complete verified runtime, while the original equations---not
the learned model or router---determine which result is returned?}
\end{quote}

This question differs from conventional and numerical solver selection~\cite{rice1976,
kotthoff2014,jessup2016lighthouse} in two ways.
First, the selected object is not a single
solver but a pipeline consisting of candidate generation, optional correction,
a separately evaluated family-specific gate, and numerical continuation after
rejection. Second, selection quality is
measured by end-to-end verified cost and numerical rejection/continuation behavior, not raw kernel latency
or hindsight winner accuracy.

We answer through \system, a verification-aware Equation-Mixture-of-Experts runtime.
It lowers typed equations to a common IR and constructs top-$k$ plans over classical,
learned, and device experts with explicit candidate, corrector, verifier, and numerical
continuation roles. A candidate is returned only after an FP64 gate evaluates the
original residual, constraints, defect, or event consistency; rejection restarts the
next numerical path from the original state.

\begin{figure}[t]
\centering
\resizebox{\columnwidth}{!}{\begin{tikzpicture}[
  node distance=5mm and 7mm,
  box/.style={draw,rounded corners,align=center,minimum height=7mm,minimum width=18mm,fill=blue!5},
  gatebox/.style={draw,rounded corners,align=center,minimum height=7mm,minimum width=18mm,fill=green!10},
  failbox/.style={draw,rounded corners,align=center,minimum height=7mm,minimum width=18mm,fill=orange!12},
  arrow/.style={-{Latex[length=2mm]},thick}
]
\node[box] (input) {Typed equation\\descriptor};
\node[box,right=of input] (ir) {Hybrid DAE IR\\and fingerprint};
\node[box,right=of ir] (assess) {Equation\\assessment};
\node[box,below=of assess] (router) {Cost--risk Router\\and top-$k$ plan};
\node[box,left=of router] (experts) {Classical, learned,\\and device experts};
\node[box,left=of experts] (correct) {Family-specific\\correction};
\node[gatebox,below=of correct] (gate) {FP64 original-equation\\recomputation gate};
\node[failbox,right=of gate] (continuation) {Next numerical\\solver path};
\node[gatebox,right=of continuation] (commit) {Verified\\result return};

\draw[arrow] (input) -- (ir);
\draw[arrow] (ir) -- (assess);
\draw[arrow] (assess) -- (router);
\draw[arrow] (router) -- (experts);
\draw[arrow] (experts) -- (correct);
\draw[arrow] (correct) -- (gate);
\draw[arrow] (gate) -- node[above,font=\scriptsize]{reject} (continuation);
\draw[arrow] (gate.south) |- node[pos=.25,left,font=\scriptsize]{accept} (commit.south);
\draw[arrow] (continuation) -- (commit);
\draw[arrow,dashed] (continuation.north) |- (router.east);
\end{tikzpicture}
}
\caption{\system's verification-aware expert-fusion pipeline. Dashed control flow
returns to the plan when an expert is inapplicable; no candidate bypasses the
family-specific original-equation gate.}
\label{fig:architecture}
\end{figure}

Our contributions are:

\begin{enumerate}
  \item We formulate heterogeneous numerical acceleration as a
  \emph{verification-aware expert-selection} problem and derive a reach-weighted
  complete-cost objective for candidate cascades, correction, original-equation
  verification, transfer, rejection, and the terminal classical solve. For a fixed eligible cascade with
  order-invariant stage statistics, we apply the classical pairwise-optimal
  cost-per-acceptance ordering rule. For finite calibrated correction actions, we
  prove exact dynamic programs for independent and adjacent-transition-cost models
  under a one-budget-per-expert constraint. We also prove the interaction-aware
  decision problem NP-complete, explaining the exponential dependence, and verify
  both optimizers against exhaustive finite oracles.

  \item We define a role-constrained candidate--corrector--gate--continuation pipeline
  in which every returned result must satisfy a family-specific original-equation
  contract. This separates numerical acceptance from router and learned-model quality:
  a poor expert may waste time but cannot bypass the gate.

  \item We implement the abstraction across algebraic and dynamic solver interfaces,
  while separating interface coverage from acceleration evidence.

  \item We retain complete-cost gains and failures. Seven PDEBench-derived workloads
  span $\PdeMinimumShort\times$--$\PdeMaximumShort\times$ median speedup with all
  95\% lower bounds above one; two operator studies beat a shared hybrid control; and
  an official-code HINTS comparison uses the same equation contract. Gate-only,
  complete-path, device, transfer, and nonlinear-routing negatives remain visible.
  Final public routing is mixed: v6 reaches $\SuiteRouteRegretPrecise\times$ oracle
  regret, only $\SuiteRouteFixedImprovementPercent\%$ below fixed, whereas the frozen
  v5 replay remains $\SuiteRouteVFiveReplayVsFixed\times$ fixed.
\end{enumerate}

The results support a bounded conclusion: verified expert fusion can accelerate
qualified repeated workloads, but it is not a universal replacement for classical
solvers. Verification cost, model applicability, correction structure, and
hardware placement must be treated as first-class design variables.
